
% =========================================================
\section{Switching geometry along the period-two branch}
\label{sec:switching-geometry}
% =========================================================

We now study how the smooth period-two branch constructed in
Section~\ref{sec:flip-bifurcation} approaches the endogenous
switching hypersurface
\[
\Sigma_c
=
\left\{
x\in\mathcal U:
\|\nabla V(x)\|_\diamond=c
\right\}.
\]

A key structural feature of gradient maps appears at this point.
For an arbitrary period-two cycle of gradient descent, the gradients
at its two points are exactly opposite. Consequently, their gradient
norms are identical for every norm $\|\cdot\|_\diamond$.

This identity has two important consequences.

First, the two points of the smooth period-two orbit necessarily
reach the clipping boundary simultaneously. Thus, unlike a generic
piecewise-smooth map, the present gradient-generated switching
problem does not exhibit a splitting of the two first-contact
parameters.

Second, the common gradient magnitude along the bifurcating cycle
has a particularly simple expansion in the square-root parameter
$\varepsilon=\sqrt{\eta-\eta_f}$. This structure will yield the
quadratic separation law in Section~\ref{sec:separation-law}.

% ---------------------------------------------------------
\subsection{Exact gradient balance on a period-two cycle}
\label{subsec:exact-gradient-balance}
% ---------------------------------------------------------

We begin with an exact identity that does not rely on the local
bifurcation expansion. An equivalent balance relation appears in
the recent characterization of gradient-descent period-two orbits
by Litman \cite{Litman2026}. Since this identity is the starting
point of the switching analysis below, we include its elementary
proof for completeness.

\begin{proposition}[Exact gradient balance]
\label{prop:exact-gradient-balance}
Let $\eta>0$, and suppose that
$\{x_-,x_+\}$ is a period-two cycle of the smooth gradient map
\[
G_\eta(x)
=
x-\eta\nabla V(x),
\]
that is,
\begin{equation}
G_\eta(x_+)=x_-,
\qquad
G_\eta(x_-)=x_+,
\qquad
x_+\neq x_-.
\label{eq:two-cycle-relations}
\end{equation}
Then
\begin{equation}
\nabla V(x_+)
=
\frac{x_+-x_-}{\eta},
\qquad
\nabla V(x_-)
=
-\frac{x_+-x_-}{\eta}.
\label{eq:exact-gradient-opposition}
\end{equation}
In particular,
\begin{equation}
\nabla V(x_+)
=
-\nabla V(x_-).
\label{eq:opposite-gradients}
\end{equation}
Consequently, for every norm $\|\cdot\|_\diamond$,
\begin{equation}
\|\nabla V(x_+)\|_\diamond
=
\|\nabla V(x_-)\|_\diamond.
\label{eq:equal-gradient-norms}
\end{equation}
\end{proposition}

\begin{proof}
The first identity in \eqref{eq:two-cycle-relations} gives
\[
x_-
=
x_+
-
\eta\nabla V(x_+),
\]
and hence
\[
\nabla V(x_+)
=
\frac{x_+-x_-}{\eta}.
\]
Similarly,
\[
x_+
=
x_-
-
\eta\nabla V(x_-)
\]
implies
\[
\nabla V(x_-)
=
\frac{x_--x_+}{\eta}
=
-\frac{x_+-x_-}{\eta}.
\]
This proves \eqref{eq:exact-gradient-opposition} and
\eqref{eq:opposite-gradients}.

Since every norm is even,
\[
\|-g\|_\diamond=\|g\|_\diamond,
\]
and therefore
\[
\|\nabla V(x_-)\|_\diamond
=
\|-\nabla V(x_+)\|_\diamond
=
\|\nabla V(x_+)\|_\diamond.
\]
\end{proof}

\begin{remark}[Gradient structure versus a generic switching map]
\label{rem:gradient-special-structure}
Equation \eqref{eq:equal-gradient-norms} is a special consequence of
the gradient-descent structure.  In a generic piecewise-smooth map,
the two points of a period-two orbit need not have equal switching
functions and may therefore encounter a switching boundary at
different parameter values.

For radial gradient clipping, by contrast, the switching function
depends only on
\[
\|\nabla V(x)\|_\diamond,
\]
and Proposition~\ref{prop:exact-gradient-balance} forces the two
switching values to be identical along every smooth period-two
cycle.
\end{remark}

% ---------------------------------------------------------
\subsection{Midpoint and half-difference coordinates}
\label{subsec:midpoint-difference}
% ---------------------------------------------------------

For the bifurcating period-two cycle
\[
\mathcal O_2(\mu)
=
\{x_-(\mu),x_+(\mu)\},
\qquad
\mu=\eta-\eta_f>0,
\]
define its midpoint and half-difference by
\begin{align}
m(\mu)
&:=
\frac{x_+(\mu)+x_-(\mu)}{2}-x_*,
\label{eq:midpoint-definition}
\\
d(\mu)
&:=
\frac{x_+(\mu)-x_-(\mu)}{2}.
\label{eq:half-difference-definition}
\end{align}
Thus
\begin{equation}
x_\pm(\mu)
=
x_*+m(\mu)\pm d(\mu).
\label{eq:midpoint-difference-representation}
\end{equation}

Theorem~\ref{thm:period-two-branch} implies
\begin{equation}
m(\mu)=O(\mu),
\qquad
d(\mu)
=
A_*\sqrt{\mu}\,v_*
+
O(\mu).
\label{eq:rough-md-orders}
\end{equation}

The period-two equations yield considerably more structure.

\begin{lemma}[Midpoint--difference equations]
\label{lem:midpoint-difference-equations}
For every smooth period-two cycle of $G_\eta$ written in the form
\eqref{eq:midpoint-difference-representation},
\begin{align}
&
\nabla V(x_*+m+d)
+
\nabla V(x_*+m-d)
=
0,
\label{eq:midpoint-equation}
\\
&
\nabla V(x_*+m+d)
-
\nabla V(x_*+m-d)
=
\frac{4}{\eta}d.
\label{eq:difference-equation}
\end{align}
\end{lemma}

\begin{proof}
By Proposition~\ref{prop:exact-gradient-balance},
\[
\nabla V(x_*+m+d)
=
\frac{2}{\eta}d
\]
and
\[
\nabla V(x_*+m-d)
=
-\frac{2}{\eta}d.
\]
Adding and subtracting these identities gives
\eqref{eq:midpoint-equation} and
\eqref{eq:difference-equation}.
\end{proof}

% ---------------------------------------------------------
\subsection{Square-root parametrization and parity}
\label{subsec:square-root-parity}
% ---------------------------------------------------------

By the parameter-dependent flip theorem
\cite{Kuznetsov2004,GuckenheimerHolmes1983} and \textbf{(H1)}, the
local period-two branch admits a $C^4$ signed parametrization in the
variable $\varepsilon$ introduced below; in particular the midpoint
and half-difference functions are $C^4$, and the parity relations
derived next are justified to the orders used in
Proposition~\ref{prop:refined-cycle-expansion}.

Introduce the signed square-root parameter
\begin{equation}
\varepsilon^2=\mu,
\qquad
\eta=\eta_f+\varepsilon^2.
\label{eq:epsilon-definition}
\end{equation}

The two-cycle equations are invariant under the transformation
\[
(d,\varepsilon)
\longmapsto
(-d,-\varepsilon),
\]
which merely exchanges the two points of the cycle.

The local uniqueness of the bifurcating period-two branch therefore
allows the branch to be parametrized so that
\begin{equation}
m(-\varepsilon)=m(\varepsilon),
\qquad
d(-\varepsilon)=-d(\varepsilon).
\label{eq:parity-md}
\end{equation}

Hence the midpoint is even and the half-difference is odd in
$\varepsilon$.

\begin{proposition}[Refined two-cycle expansion]
\label{prop:refined-cycle-expansion}
Assume \textbf{(H1)--(H4)}.
The period-two branch can be parametrized by the signed variable
$\varepsilon$, with $\mu=\varepsilon^2$, such that
\begin{align}
m(\varepsilon)
&=
\varepsilon^2m_2
+
O(\varepsilon^4),
\label{eq:m-expansion}
\\
d(\varepsilon)
&=
A_*\varepsilon v_*
+
O(\varepsilon^3),
\label{eq:d-expansion}
\end{align}
where
\begin{equation}
m_2
=
-\frac{A_*^2}{2}H_*^{-1}g_*
=
-\frac{A_*^2}{2}h_*.
\label{eq:m2-formula}
\end{equation}

Consequently,
\begin{equation}
x_\pm(\mu)
=
x_*
-
\frac{A_*^2}{2}\mu h_*
\pm
A_*\sqrt{\mu}\,v_*
+
O(\mu^{3/2}).
\label{eq:refined-xpm-expansion}
\end{equation}
\end{proposition}

\begin{proof}
The parity relations \eqref{eq:parity-md} imply that the local
expansions have the form
\[
m(\varepsilon)
=
\varepsilon^2m_2+O(\varepsilon^4)
\]
and
\[
d(\varepsilon)
=
\varepsilon d_1+O(\varepsilon^3).
\]
By Theorem~\ref{thm:period-two-branch},
\[
d_1=A_*v_*.
\]

It remains to determine $m_2$.
Taylor expansion of the gradient about $x_*$ gives
\begin{equation}
\nabla V(x_*+y)
=
H_*y
+
\frac12
T_*[y,y,\cdot]^\sharp
+
O(\|y\|_2^3).
\label{eq:gradient-second-order-Taylor}
\end{equation}

Apply \eqref{eq:gradient-second-order-Taylor} to
$y=m+d$ and $y=m-d$ and add the resulting expressions.
Using \eqref{eq:midpoint-equation}, we obtain
\begin{equation}
2H_*m
+
T_*[d,d,\cdot]^\sharp
+
O(\varepsilon^4)
=
0.
\label{eq:midpoint-expanded}
\end{equation}
Here we used
\[
m=O(\varepsilon^2),
\qquad
d=O(\varepsilon),
\]
and the cancellation of all terms odd in $d$.

Substituting
\[
m=\varepsilon^2m_2+O(\varepsilon^4)
\]
and
\[
d=A_*\varepsilon v_*+O(\varepsilon^3)
\]
into \eqref{eq:midpoint-expanded} gives
\[
2H_*m_2
+
A_*^2
T_*[v_*,v_*,\cdot]^\sharp
=
0.
\]
Since $H_*$ is invertible,
\[
m_2
=
-\frac{A_*^2}{2}
H_*^{-1}
T_*[v_*,v_*,\cdot]^\sharp
=
-\frac{A_*^2}{2}h_*.
\]
This proves \eqref{eq:m2-formula} and hence
\eqref{eq:refined-xpm-expansion}.
\end{proof}

\begin{remark}[Nonlinear displacement of the cycle midpoint]
\label{rem:midpoint-drift}
Although the leading oscillation is aligned with the maximal
Hessian eigendirection $v_*$, the midpoint of the bifurcating
two-cycle generally does not remain at $x_*$. Instead,
\[
\frac{x_+(\mu)+x_-(\mu)}{2}
=
x_*
-
\frac{A_*^2}{2}\mu h_*
+
O(\mu^2).
\]
Thus the third derivative of $V$ produces an $O(\mu)$ drift of the
cycle center through the vector
\[
h_*
=
H_*^{-1}
T_*[v_*,v_*,\cdot]^\sharp.
\]
\end{remark}

% ---------------------------------------------------------
\subsection{A consistency derivation of the flip amplitude}
\label{subsec:amplitude-consistency}
% ---------------------------------------------------------

The midpoint--difference equations also provide a direct,
coordinate-free derivation of the leading oscillation amplitude.

\begin{proposition}[Amplitude identity]
\label{prop:amplitude-identity}
Under \textbf{(H1)--(H4)}, the leading amplitude satisfies
\begin{equation}
A_*^2\Gamma_*
=
3\lambda_*^2.
\label{eq:amplitude-Gamma-identity}
\end{equation}
Equivalently,
\begin{equation}
A_*^2
=
\frac{3\lambda_*^2}{\Gamma_*},
\label{eq:amplitude-coordinate-free}
\end{equation}
in agreement with
\eqref{eq:A-star-Gamma}.
\end{proposition}

\begin{proof}
Expand the difference equation
\eqref{eq:difference-equation} to cubic order.

Using
\[
m
=
\varepsilon^2m_2+O(\varepsilon^4),
\qquad
d
=
A_*\varepsilon v_*+O(\varepsilon^3),
\]
Taylor expansion gives
\begin{align}
&
\nabla V(x_*+m+d)
-
\nabla V(x_*+m-d)
\nonumber\\
&\quad=
2H_*d
+
2T_*[m,d,\cdot]^\sharp
+
\frac13Q_*[d,d,d,\cdot]^\sharp
+
O(\varepsilon^4).
\label{eq:difference-expanded}
\end{align}

On the other hand,
\[
\frac{4}{\eta}
=
\frac{4}{\eta_f+\varepsilon^2}
=
\frac{4}{\eta_f}
-
\frac{4}{\eta_f^2}\varepsilon^2
+
O(\varepsilon^4).
\label{eq:four-over-eta}
\]

Take the Euclidean inner product of
\eqref{eq:difference-equation} with $v_*$ and compare the
coefficients of $\varepsilon^3$.

The unknown $O(\varepsilon^3)$ correction in $d$ cancels because
\[
\frac{4}{\eta_f}
=
2\lambda_*.
\]
Thus,
\begin{equation}
2A_*
T_*[v_*,v_*,m_2]
+
\frac{A_*^3}{3}
Q_*[v_*,v_*,v_*,v_*]
=
-\frac{4A_*}{\eta_f^2}.
\label{eq:amplitude-before-substitution}
\end{equation}

Since
\[
m_2
=
-\frac{A_*^2}{2}h_*
\]
and
\[
\eta_f=\frac{2}{\lambda_*},
\]
equation \eqref{eq:amplitude-before-substitution} becomes
\[
-A_*^3T_*[v_*,v_*,h_*]
+
\frac{A_*^3}{3}Q_*[v_*^4]
=
-A_*\lambda_*^2.
\]
Dividing by $A_*>0$ and rearranging gives
\[
A_*^2
\left(
3T_*[v_*,v_*,h_*]
-
Q_*[v_*^4]
\right)
=
3\lambda_*^2.
\]
By definition of $\Gamma_*$,
\[
A_*^2\Gamma_*
=
3\lambda_*^2.
\]
\end{proof}

\begin{remark}
\label{rem:coordinate-free-check}
Proposition~\ref{prop:amplitude-identity} provides an independent
check of the center-manifold normal-form calculation in
Section~\ref{sec:flip-bifurcation}.  In particular, the same
high-order quantity $\Gamma_*$ governs both the supercriticality
criterion and the leading amplitude of the emerging period-two
cycle.
\end{remark}

% ---------------------------------------------------------
\subsection{Gradient exposure along the bifurcating cycle}
\label{subsec:gradient-exposure}
% ---------------------------------------------------------

Define the common gradient exposure of the smooth period-two cycle
by
\begin{equation}
\mathscr E(\mu)
:=
\|\nabla V(x_+(\mu))\|_\diamond
=
\|\nabla V(x_-(\mu))\|_\diamond.
\label{eq:exposure-definition}
\end{equation}
The equality is exact by
Proposition~\ref{prop:exact-gradient-balance}.

Using the half-difference coordinate, one obtains the exact formula
\begin{equation}
\mathscr E(\mu)
=
\frac{2}{\eta_f+\mu}
\|d(\mu)\|_\diamond.
\label{eq:exact-exposure-d}
\end{equation}

This identity is considerably stronger than a direct Taylor
expansion of $\nabla V$ at the two cycle points.

\begin{theorem}[Asymptotic switching geometry of the two-cycle]
\label{thm:exposure-expansion}
Assume \textbf{(H1)--(H5)}.
As $\mu\to0^+$,
\begin{equation}
\mathscr E(\mu)
=
B_*\mu^{1/2}
+
O(\mu^{3/2}),
\label{eq:exposure-expansion}
\end{equation}
where
\begin{equation}
B_*
=
\lambda_*A_*
\|v_*\|_\diamond.
\label{eq:B-star}
\end{equation}

Equivalently,
\begin{equation}
B_*^2
=
\frac{
3\lambda_*^4
\|v_*\|_\diamond^2
}{
\Gamma_*
}.
\label{eq:B-star-Gamma}
\end{equation}

In particular, there is no term of order $\mu$ in
\eqref{eq:exposure-expansion}.
\end{theorem}

\begin{proof}
Set
\[
\varepsilon=\sqrt{\mu}.
\]
By Proposition~\ref{prop:refined-cycle-expansion},
\[
d(\varepsilon)
=
A_*\varepsilon v_*
+
O(\varepsilon^3).
\]
Using positive homogeneity of the norm,
\begin{align}
\|d(\varepsilon)\|_\diamond
&=
\varepsilon
\left\|
A_*v_*
+
O(\varepsilon^2)
\right\|_\diamond
\nonumber\\
&=
A_*
\|v_*\|_\diamond
\varepsilon
+
O(\varepsilon^3).
\label{eq:d-norm-expansion}
\end{align}
Furthermore,
\begin{equation}
\frac{2}{\eta_f+\varepsilon^2}
=
\frac{2}{\eta_f}
+
O(\varepsilon^2)
=
\lambda_*
+
O(\varepsilon^2).
\label{eq:two-over-eta-expansion}
\end{equation}
Substituting
\eqref{eq:d-norm-expansion} and
\eqref{eq:two-over-eta-expansion} into
\eqref{eq:exact-exposure-d} yields
\[
\mathscr E(\mu)
=
\lambda_*A_*
\|v_*\|_\diamond
\varepsilon
+
O(\varepsilon^3).
\]
Since $\varepsilon=\sqrt{\mu}$,
\[
\mathscr E(\mu)
=
B_*\mu^{1/2}
+
O(\mu^{3/2}),
\]
where
\[
B_*
=
\lambda_*A_*
\|v_*\|_\diamond.
\]

Finally, using
\[
A_*^2
=
\frac{3\lambda_*^2}{\Gamma_*}
\]
gives
\[
B_*^2
=
\lambda_*^2
A_*^2
\|v_*\|_\diamond^2
=
\frac{
3\lambda_*^4
\|v_*\|_\diamond^2
}{
\Gamma_*
}.
\]
\end{proof}

\begin{remark}[Absence of a quadratic term in the exposure]
\label{rem:no-mu-term}
A naive substitution of
\[
x_\pm(\mu)
=
x_*
\pm
A_*\sqrt{\mu}\,v_*
+
O(\mu)
\]
into the gradient norm might suggest an expansion of the form
\[
B_1\sqrt{\mu}+B_2\mu+\cdots.
\]
The exact two-cycle identity
\eqref{eq:exact-exposure-d} shows that this is not the case.

The half-difference $d$ is odd in the signed square-root parameter
$\varepsilon$, while $\eta=\eta_f+\varepsilon^2$ is even.
Consequently, the common exposure is an odd function of
$\varepsilon$ to leading orders:
\[
\mathscr E(\varepsilon^2)
=
B_*\varepsilon
+
O(\varepsilon^3).
\]
This structural cancellation will improve the remainder in the
switching-contact separation law from the naively expected
$O(c^3)$ to $O(c^4)$.
\end{remark}

% ---------------------------------------------------------
\subsection{Simultaneous first contact}
\label{subsec:simultaneous-contact}
% ---------------------------------------------------------

We now record the most important geometric consequence of the exact
gradient balance.

\begin{theorem}[Simultaneous switching contact]
\label{thm:simultaneous-contact}
Let $\mathcal O_2(\mu)=\{x_-(\mu),x_+(\mu)\}$ be the smooth
period-two branch from
Theorem~\ref{thm:period-two-branch}.
For every sufficiently small $\mu>0$,
\begin{equation}
x_+(\mu)\in\Sigma_c
\quad\Longleftrightarrow\quad
x_-(\mu)\in\Sigma_c.
\label{eq:simultaneous-contact-equivalence}
\end{equation}

Likewise,
\begin{align}
x_+(\mu)\in\Omega_c^-
&\quad\Longleftrightarrow\quad
x_-(\mu)\in\Omega_c^-,
\label{eq:simultaneous-smooth}
\\
x_+(\mu)\in\Omega_c^+
&\quad\Longleftrightarrow\quad
x_-(\mu)\in\Omega_c^+.
\label{eq:simultaneous-clipped}
\end{align}

Hence the two points of a smooth gradient-descent period-two cycle
cannot cross a radial clipping boundary sequentially.
\end{theorem}

\begin{proof}
By Proposition~\ref{prop:exact-gradient-balance},
\[
\|\nabla V(x_+(\mu))\|_\diamond
=
\|\nabla V(x_-(\mu))\|_\diamond.
\]
Each of the three statements follows immediately by comparing this
common value with the clipping threshold $c$.
\end{proof}

\begin{corollary}[No pre-contact splitting of the two-cycle]
\label{cor:no-contact-splitting}
There do not exist two distinct first-contact parameters
\[
\eta_{\mathrm{sc}}^+(c)
\neq
\eta_{\mathrm{sc}}^-(c)
\]
for the two points of the smooth period-two branch.
Whenever the first switching contact exists, it is a simultaneous
two-point contact.
\end{corollary}

\begin{proof}
If one point lies on $\Sigma_c$, then the other lies on $\Sigma_c$
by Theorem~\ref{thm:simultaneous-contact}.
\end{proof}

% ---------------------------------------------------------
\subsection{Impossibility of a mixed smooth--clipped two-cycle}
\label{subsec:no-mixed-two-cycle}
% ---------------------------------------------------------

The simultaneous-contact structure is not restricted to the smooth
branch before contact.  Radial clipping itself rules out a genuine
period-two cycle having one point in the smooth region and the other
in the clipped region.

\begin{proposition}[No mixed-itinerary period-two cycle]
\label{prop:no-mixed-period-two}
Let
\[
\{y_-,y_+\}
\]
be a period-two cycle of the clipped map
\[
F_{\eta,c}(x)
=
x-\eta
\mathcal C_c^\diamond(\nabla V(x)).
\]
Then it is impossible to have
\begin{equation}
y_-\in\Omega_c^-,
\qquad
y_+\in\Omega_c^+,
\label{eq:forbidden-mixed-cycle}
\end{equation}
or vice versa.

Therefore every period-two cycle of $F_{\eta,c}$ is either entirely
smooth, entirely clipped, or lies on the switching boundary.
\end{proposition}

\begin{proof}
The period-two relations give
\begin{align}
\mathcal C_c^\diamond(\nabla V(y_+))
&=
\frac{y_+-y_-}{\eta},
\label{eq:clipped-balance-plus}
\\
\mathcal C_c^\diamond(\nabla V(y_-))
&=
-\frac{y_+-y_-}{\eta}.
\label{eq:clipped-balance-minus}
\end{align}
Hence
\begin{equation}
\left\|
\mathcal C_c^\diamond(\nabla V(y_+))
\right\|_\diamond
=
\left\|
\mathcal C_c^\diamond(\nabla V(y_-))
\right\|_\diamond.
\label{eq:clipped-output-equal}
\end{equation}

For radial clipping,
\begin{equation}
\left\|
\mathcal C_c^\diamond(g)
\right\|_\diamond
=
\min\{\|g\|_\diamond,c\}.
\label{eq:clipped-output-norm}
\end{equation}

Suppose, for contradiction, that
\[
y_-\in\Omega_c^-,
\qquad
y_+\in\Omega_c^+.
\]
Then
\[
\left\|
\mathcal C_c^\diamond(\nabla V(y_-))
\right\|_\diamond
=
\|\nabla V(y_-)\|_\diamond
<
c,
\]
whereas
\[
\left\|
\mathcal C_c^\diamond(\nabla V(y_+))
\right\|_\diamond
=
c.
\]
This contradicts \eqref{eq:clipped-output-equal}.
The opposite mixed itinerary is ruled out identically.
\end{proof}

\begin{remark}[Structural consequence for the post-contact problem]
\label{rem:no-mixed-consequence}
Proposition~\ref{prop:no-mixed-period-two} rules out the
mixed smooth--clipped period-two continuation that one might expect
from a generic switching system.

Thus the post-contact problem for radial gradient clipping is more
constrained: if a period-two branch persists beyond the switching
contact, both points must enter the clipped regime together.

The corresponding post-contact continuation problem will therefore
be formulated as a simultaneous two-point switching bifurcation,
rather than as a mixed-itinerary continuation.
\end{remark}

% ---------------------------------------------------------
\subsection{Summary of the switching geometry}
\label{subsec:switching-geometry-summary}
% ---------------------------------------------------------

The preceding analysis establishes three structural facts that will
be used in the next section.

First, the period-two gradients are exactly opposite:
\begin{equation}
\nabla V(x_+(\mu))
=
-\nabla V(x_-(\mu)).
\label{eq:summary-opposite}
\end{equation}

Second, their common clipping exposure satisfies
\begin{equation}
\mathscr E(\mu)
=
\lambda_*A_*
\|v_*\|_\diamond
\sqrt{\mu}
+
O(\mu^{3/2}).
\label{eq:summary-exposure}
\end{equation}

Third, the first switching interaction is necessarily simultaneous
at both points of the cycle.

Therefore the first-contact condition reduces to the single scalar
equation
\begin{equation}
\mathscr E(\mu_{\mathrm{sc}}(c))
=
c.
\label{eq:first-contact-scalar}
\end{equation}

Section~\ref{sec:separation-law} will invert
\eqref{eq:first-contact-scalar} and derive the resulting separation
law between the smooth flip threshold $\eta_f$ and the first
switching-contact parameter.

